\documentclass[a4paper,11pt]{article}

\usepackage[italian]{babel}
\usepackage[T1]{fontenc}
\usepackage[utf8]{inputenc}
\usepackage{lmodern}
\usepackage{geometry}
\usepackage{xcolor}
\usepackage{hyperref}
\usepackage{tabularx}
\usepackage{array}
\usepackage{booktabs}
\usepackage{enumitem}
\usepackage{titlesec}
\usepackage{tcolorbox}
\usepackage{fancyhdr}
\usepackage{multicol}
\usepackage{amsmath}

\geometry{margin=2.2cm}
\setlength{\parskip}{0.6em}
\setlength{\parindent}{0pt}

% --- Colori ---
\definecolor{MainColor}{HTML}{1F4E79}
\definecolor{AccentColor}{HTML}{2E86AB}
\definecolor{SoftBlue}{HTML}{EAF4FB}
\definecolor{SoftGray}{HTML}{F5F7FA}
\definecolor{SoftGreen}{HTML}{EAF8F0}
\definecolor{DarkText}{HTML}{1E1E1E}

% --- Hyperref ---
\hypersetup{
    colorlinks=true,
    linkcolor=MainColor,
    urlcolor=AccentColor,
    pdftitle={Guida moderna a tmux},
    pdfauthor={}
}

% --- Titoli ---
\titleformat{\section}
  {\Large\bfseries\color{MainColor}}
  {\thesection}{0.8em}{}

\titleformat{\subsection}
  {\large\bfseries\color{AccentColor}}
  {\thesubsection}{0.8em}{}

% --- Header/Footer ---
\pagestyle{fancy}
\fancyhf{}
\fancyhead[L]{\textcolor{MainColor}{Guida a tmux}}
\fancyhead[R]{\textcolor{AccentColor}{Terminal multiplexer}}
\fancyfoot[C]{\thepage}

% --- Box ---
\newtcolorbox{infobox}{
    colback=SoftBlue,
    colframe=AccentColor,
    boxrule=0.8pt,
    arc=3mm,
    left=2mm,right=2mm,top=2mm,bottom=2mm
}

\newtcolorbox{tipbox}{
    colback=SoftGreen,
    colframe=green!50!black,
    boxrule=0.8pt,
    arc=3mm,
    left=2mm,right=2mm,top=2mm,bottom=2mm
}

\newtcolorbox{codebox}{
    colback=SoftGray,
    colframe=black!20,
    boxrule=0.5pt,
    arc=2mm,
    left=2mm,right=2mm,top=2mm,bottom=2mm
}

% --- Comando rapido ---
\newcommand{\shortcut}[2]{\texttt{#1} & #2 \\}

\begin{document}

% =========================
% Titolo
% =========================
\vspace*{1.5cm}

{\Huge\bfseries\color{MainColor} Guida  a \texttt{tmux}} \\[0.4cm]
{\Large\color{AccentColor} Per iniziare e per avere sempre i comandi a portata di mano}


\vspace{1cm}

\begin{infobox}
\textbf \texttt{tmux} ti permette di aprire più terminali, dividerli in pannelli, scollegarti senza fermare i processi e riprendere tutto in un secondo momento.
\end{infobox}

\vspace{0.5cm}

\begin{codebox}
\textbf{Quick reference}

\renewcommand{\arraystretch}{1.2}
\begin{tabularx}{\textwidth}{>{\ttfamily}l X}
\toprule
Comando & Descrizione \\
\midrule
 tmux new -s nome & Creazione sessione \\
 tmux ls & Lista sessioni \\
 tmux attach -t nome & Attach a sessione \\
 Ctrl+b d & Detach sessione \\
 Ctrl+b c & Nuova finestra \\
 Ctrl+b \% & Split verticale \\
 Ctrl+b " & Split orizzontale \\
 Ctrl+b frecce & Navigazione pannelli \\
 Ctrl+b x & Chiusura pannello \\
 Ctrl+b ? & Help \\
\bottomrule
\end{tabularx}
\end{codebox}

\vspace{0.6cm}
\tableofcontents
\newpage

% =========================
\section{Cos'è tmux}
\texttt{tmux} è un \textbf{terminal multiplexer}, cioè uno strumento che ti consente di lavorare con più shell dentro una sola sessione di terminale.

È particolarmente utile quando:
\begin{itemize}[leftmargin=1.2cm]
    \item devi lanciare processi lunghi;
    \item lavori su server remoti tramite SSH;
    \item vuoi dividere il terminale in più aree;
    \item vuoi chiudere il terminale senza interrompere ciò che sta girando.
\end{itemize}

\begin{tipbox}
\textbf{Idea pratica:} puoi avviare un training, uno script o una compilazione dentro \texttt{tmux}, scollegarti, e poi rientrare più tardi trovando tutto ancora attivo.
\end{tipbox}

% =========================
\section{I tre concetti fondamentali}
Per usare bene \texttt{tmux}, basta capire tre oggetti:

\begin{infobox}
\textbf{Sessione} \\
È il contenitore principale. Puoi pensarla come un workspace completo.

\medskip
\textbf{Finestra (window)} \\
È simile a una scheda del browser: una sessione può avere più finestre.

\medskip
\textbf{Pannello (pane)} \\
È una divisione della finestra. Una finestra può essere spezzata in più parti.
\end{infobox}

Quindi la struttura mentale è:

\[
\text{Sessione} \rightarrow \text{Finestre} \rightarrow \text{Pannelli}
\]

% =========================
\section{Primo utilizzo}
\subsection{Avviare tmux}
\begin{codebox}
\texttt{tmux}
\end{codebox}

Questo apre una nuova sessione.

\subsection{Creare una sessione con nome}
\begin{codebox}
\texttt{tmux new -s lavoro}
\end{codebox}

Qui \texttt{lavoro} è il nome della sessione. Dare un nome è molto utile, soprattutto se lavori su più progetti.

\subsection{Vedere le sessioni attive}
\begin{codebox}
\texttt{tmux ls}
\end{codebox}

\subsection{Rientrare in una sessione}
\begin{codebox}
\texttt{tmux attach -t lavoro}
\end{codebox}

% =========================
\section{Il prefisso dei comandi}
Quasi tutti gli shortcut di \texttt{tmux} iniziano con un \textbf{prefisso}. Di default il prefisso è:

\begin{center}
\Large \texttt{Ctrl+b}
\end{center}

Questo significa che molti comandi si usano così:
\begin{itemize}[leftmargin=1.2cm]
    \item premi \texttt{Ctrl+b}
    \item rilascia
    \item premi il tasto del comando
\end{itemize}

Per esempio:
\begin{codebox}
\texttt{Ctrl+b} poi \texttt{c}
\end{codebox}
crea una nuova finestra.

% =========================
\section{Comandi essenziali}
\subsection{Staccarsi dalla sessione senza chiuderla}
\begin{codebox}
\texttt{Ctrl+b} poi \texttt{d}
\end{codebox}

Questo comando fa il \textbf{detach}: esci dalla sessione ma i processi continuano a girare.

\subsection{Chiudere davvero un pannello o una finestra}
Per chiudere un pannello o una finestra, in genere basta uscire dalla shell:
\begin{codebox}
\texttt{exit}
\end{codebox}

oppure usare:
\begin{codebox}
\texttt{Ctrl+d}
\end{codebox}

% =========================
\section{Finestre}
Le finestre servono a separare attività diverse: per esempio editor, log, training, shell principale.

\begin{center}
\renewcommand{\arraystretch}{1.3}
\begin{tabularx}{\textwidth}{>{\ttfamily}l X}
\toprule
Comando & Funzione \\
\midrule
Ctrl+b c & Crea una nuova finestra \\
Ctrl+b n & Va alla finestra successiva \\
Ctrl+b p & Va alla finestra precedente \\
Ctrl+b w & Mostra la lista delle finestre \\
Ctrl+b , & Rinomina la finestra corrente \\
Ctrl+b 0...9 & Va direttamente alla finestra con quel numero \\
\bottomrule
\end{tabularx}
\end{center}

\begin{tipbox}
\textbf{Consiglio utile:} rinomina le finestre con nomi come \texttt{editor}, \texttt{logs}, \texttt{train}, \texttt{server}. Ti orienti molto meglio.
\end{tipbox}

% =========================
\section{Pannelli}
I pannelli sono molto comodi quando vuoi vedere più terminali contemporaneamente nella stessa finestra.

\subsection{Dividere lo schermo}
\begin{center}
\renewcommand{\arraystretch}{1.3}
\begin{tabularx}{\textwidth}{>{\ttfamily}l X}
\toprule
Comando & Funzione \\
\midrule
Ctrl+b \% & Divide verticalmente \\
Ctrl+b " & Divide orizzontalmente \\
Ctrl+b frecce & Passa al pannello vicino \\
Ctrl+b x & Chiude il pannello corrente \\
Ctrl+b z & Ingrandisce/riduce il pannello corrente \\
\bottomrule
\end{tabularx}
\end{center}

\subsection{Idea pratica}
Una disposizione classica può essere:
\begin{itemize}[leftmargin=1.2cm]
    \item pannello sinistro: editor o shell principale;
    \item pannello in alto a destra: esecuzione script;
    \item pannello in basso a destra: log o monitoraggio.
\end{itemize}

% =========================
\section{Sessioni}
Le sessioni sono la vera forza di \texttt{tmux}.

\begin{center}
\renewcommand{\arraystretch}{1.3}
\begin{tabularx}{\textwidth}{>{\ttfamily}l X}
\toprule
Comando & Funzione \\
\midrule
tmux new -s nome & Crea una nuova sessione con nome \\
tmux ls & Elenca le sessioni attive \\
tmux attach -t nome & Si collega a una sessione esistente \\
Ctrl+b d & Si scollega dalla sessione senza fermarla \\
Ctrl+b s & Mostra le sessioni disponibili \\
tmux kill-session -t nome & Chiude completamente una sessione \\
\bottomrule
\end{tabularx}
\end{center}

\begin{infobox}
\textbf{Detach} non significa chiudere. \\
Significa solo uscire temporaneamente dalla sessione lasciandola viva.
\end{infobox}

% =========================
\section{Workflow tipico}
Questa è una sequenza molto comune nell'uso reale.

\subsection{Caso: lanciare un processo lungo}
\begin{enumerate}[leftmargin=1.2cm]
    \item Crea una sessione:
\begin{codebox}
\texttt{tmux new -s esperimento}
\end{codebox}

    \item Avvia il comando che durerà a lungo:
\begin{codebox}
\texttt{python train.py}
\end{codebox}

    \item Staccati:
\begin{codebox}
\texttt{Ctrl+b} poi \texttt{d}
\end{codebox}

    \item Più tardi rientra:
\begin{codebox}
\texttt{tmux attach -t esperimento}
\end{codebox}
\end{enumerate}

Questo è uno degli usi più importanti di \texttt{tmux}.

% =========================
\section{Cheat sheet veloce}
Questa sezione è pensata per chi vuole ricordare rapidamente i comandi principali.

\begin{center}
\renewcommand{\arraystretch}{1.25}
\begin{tabularx}{\textwidth}{>{\ttfamily}l X}
\toprule
Shortcut / comando & Significato \\
\midrule
tmux & Avvia una sessione nuova \\
tmux new -s nome & Crea una sessione con nome \\
tmux ls & Elenca le sessioni \\
tmux attach -t nome & Rientra in una sessione \\
Ctrl+b d & Detach dalla sessione \\
Ctrl+b c & Nuova finestra \\
Ctrl+b n & Finestra successiva \\
Ctrl+b p & Finestra precedente \\
Ctrl+b \% & Split verticale \\
Ctrl+b " & Split orizzontale \\
Ctrl+b frecce & Muoversi tra pannelli \\
Ctrl+b x & Chiudere il pannello \\
Ctrl+b z & Zoom del pannello \\
Ctrl+b w & Lista finestre \\
Ctrl+b s & Lista sessioni \\
Ctrl+b ? & Aiuto interno di tmux \\
\bottomrule
\end{tabularx}
\end{center}

% =========================
\section{Comandi utili da ricordare}
\subsection{Aiuto interno}
\begin{codebox}
\texttt{Ctrl+b} poi \texttt{?}
\end{codebox}

Mostra la lista dei comandi disponibili dentro \texttt{tmux}.

\subsection{Modalità copia}
In molte configurazioni, \texttt{tmux} supporta una modalità di scroll e copia. Anche se all'inizio può non servirti, è bene sapere che esiste.

\subsection{Mouse}
Puoi abilitare il mouse nella configurazione per:
\begin{itemize}[leftmargin=1.2cm]
    \item cambiare pannello con un clic;
    \item ridimensionare pannelli;
    \item scorrere meglio.
\end{itemize}

% =========================
\section{Configurazione minima consigliata}
Nel file:
\begin{codebox}
\texttt{\~/.tmux.conf}
\end{codebox}

puoi aggiungere una configurazione base come questa:

\begin{codebox}
\texttt{set -g mouse on} \\
\texttt{set -g history-limit 10000}
\end{codebox}

\textbf{Cosa fanno:}
\begin{itemize}[leftmargin=1.2cm]
    \item \texttt{mouse on}: abilita il supporto al mouse;
    \item \texttt{history-limit 10000}: aumenta la cronologia del terminale.
\end{itemize}

% =========================
\section{Errori comuni dei principianti}
\begin{infobox}
\textbf{1. Pensare che chiudere il terminale chiuda tutto} \\
No: se il processo è dentro una sessione \texttt{tmux} attiva, puoi rientrare dopo.

\medskip
\textbf{2. Confondere sessione, finestra e pannello} \\
Sessione = contenitore grande; finestra = tab; pannello = divisione.

\medskip
\textbf{3. Dimenticare il prefisso} \\
Se uno shortcut non funziona, probabilmente manca prima \texttt{Ctrl+b}.
\end{infobox}

% =========================
\section{Quando usare tmux}
\begin{itemize}[leftmargin=1.2cm]
    \item quando lavori su server remoti;
    \item quando lanci script lunghi;
    \item quando vuoi più terminali nello stesso spazio;
    \item quando vuoi organizzare meglio il lavoro da shell.
\end{itemize}

% =========================
\section{Conclusione}
\texttt{tmux} all'inizio può sembrare un po' strano, ma dopo poche ore diventa uno strumento estremamente naturale.

\begin{tipbox}
\textbf{Da ricordare:}
\begin{itemize}[leftmargin=1.2cm]
    \item crea una sessione con nome;
    \item usa finestre per separare i compiti;
    \item usa pannelli per vedere più terminali insieme;
    \item fai detach quando vuoi uscire senza interrompere nulla.
\end{itemize}
\end{tipbox}

\vfill
\begin{center}
\textcolor{AccentColor}{\large Buon lavoro con \texttt{tmux}}
\end{center}

\end{document}